
% Question 11
\question
An inverted right circular conical vessel contains some milk. The vessel is held vertically. The depth of milk in the vessel is $12\text{ cm}$. Peter then pours $444\pi\text{ cm}^3$ of milk into the vessel without overflowing. He now finds that the depth of milk in the vessel is $16\text{ cm}$.
\begin{parts}
\part Express the final volume of milk in the vessel in terms of $\pi$. \rmarks{3}
\part Peter claims that the final area of the wet curved surface of the vessel is at least $800\text{ cm}^2$. Do you agree? Explain your answer. \rmarks{3}
\end{parts}

\begin{solution}
\begin{parts}
\part Let $V\text{ cm}^3$ be the final volume of milk in the vessel.
\begin{align*}
\frac{V - 444\pi}{V} &= \left(\frac{12}{16}\right)^3 \quad &\text{(1M + 1A)} \\
V &= 768\pi \quad &\text{(1A)}
\end{align*}
Thus, the final volume of milk in the vessel is $768\pi\text{ cm}^3$.
\part Let $r\text{ cm}$ be the final radius of the surface of milk in the vessel.
\begin{align*}
\frac{1}{3}\pi r^2(16) &= 768\pi \quad &\text{(1M)} \\
r &= 12
\end{align*}
The final area of the wet curved surface of the vessel is
\begin{align*}
&= \pi(12)\sqrt{12^2 + 16^2} \quad &\text{(1M)} \\
&= 240\pi \approx 753.9822369\text{ cm}^2 < 800\text{ cm}^2
\end{align*}
Thus, the claim is disagreed. \hfill (1A f.t.)
\end{parts}
\end{solution}
